% Auto-generado por scripts/extract-lessons-to-tex.ps1
% Fuente: HTML en carpeta L1-Fundamentos-ontologicos-convencionalidad-y-limite-al-ius-puniendi (sin modificar archivos originales).

\documentclass[11pt,a4paper]{article}
\usepackage[utf8]{inputenc}
\usepackage[T1]{fontenc}
\usepackage[spanish]{babel}
\usepackage{geometry}
\geometry{margin=2.5cm}
\usepackage{hyperref}
\hypersetup{colorlinks=true, linkcolor=blue, urlcolor=blue}

\title{\texttt{L1-Fundamentos-ontologicos-convencionalidad-y-limite-al-ius-puniendi}}
\author{Extracci\'on desde lecciones HTML}
\date{}

\begin{document}
\maketitle

\begin{small}
\noindent Este documento consolida el texto visible de los ocho componentes (01--08). Los formularios interactivos aparecen como texto; la l\'ogica JavaScript no se reproduce aqu\'i.
\end{small}
\bigskip
\section{Conceptos}
\noindent\textit{Archivo fuente:} \texttt{01-conceptos.html}

L1 — Conceptos\par





    Elemento 1 — Conceptos\par
    Teoría de los principios y normas rectoras\par

    Introducción breve: aquí se fijan las categorías dogmáticas que condicionan todo el Capítulo 4 de Bernal.\par


      Este elemento sistematiza la teleología constitucional del derecho penal, el bloque de constitucionalidad y la\par
      legalidad estricta como límite material y formal al ius puniendi. La lectura debe hacerse en clave de control:\par
      qué principios invalidan o reinterpretan normas penales y actuaciones procesales.\par



      1. Principios rectores y función sistémica\par


        Los principios rectores no son ornamentos retóricos: articulan el mínimo constitucional que debe respetar el legislador, el\par
        intérprete y el aplicador del derecho penal. Su violación grave puede proyectarse en inconstitucionalidad del tipo, nulidades o líneas\par
        defensivas en instancias superiores.\par



        1.1 Jerarquía argumentativa\par

          En litigio complejo, la secuencia típica es: (i) identificar el estándar constitucional o convencional aplicable; (ii) contrastar la\par
          norma o la decisión con ese estándar; (iii) deducir el remedio procesal idóneo. Los principios rectores orientan el paso (i) y el\par
          estándar de revisión del paso (ii).\par




        1.2 Relación con el Código Penal colombiano\par

          Los arts. 1 a 6 del C.P. expresan parte de esa arquitectura (fines del código, principios de las sanciones, funciones de la pena).\par
          El dogmatista debe leerlos junto con la Constitución y el bloque de constitucionalidad, no como texto aislado.\par






      2. Dignidad humana como principio teleológico\par


        La dignidad actúa como criterio de proporcionalidad y humanización de la pena y del proceso. Impide tratar al imputado\par
        como mero objeto de política criminal y exige que la severidad y el procedimiento sean compatibles con la persona en su singularidad.\par



        Referencia de profundización: Bernal, Cap. 4, sección dedicada al principio de dignidad humana (índice del libro, p. 266).\par





      3. Bloque de constitucionalidad: lato sensu y stricto sensu\par



          Lato sensu\par


            Comprende el conjunto amplio de fuentes y estándares que informan la interpretación conforme a la Constitución: texto constitucional,\par
            leyes estatutarias, tratados y, según el caso, desarrollo jurisprudencial que integra derechos y deberes estatales.\par





          Stricto sensu\par


            Alude al núcleo duro de supremacía: normas y controles que vinculan directamente al juez penal y al legislador cuando hay colisión\par
            con el texto constitucional o estándares convencionales incorporados.\par







        En práctica forense, la distinción ayuda a decidir qué argumento va primero: integración interpretativa amplia frente a control\par
        concentrado o excepciones de constitucionalidad, según la vía procesal.\par





      4. Legalidad estricta: lex certa, praevia, stricta\par


        La tríada legal es el límite formal al delito y a la pena. Opera en estrecha conexión con la tipicidad y con el derecho\par
        de defensa: sin tipo predeterminado y claro, la defensa técnica se ve comprometida materialmente.\par




          Lex certa\par

            Exige precisión del tipo: conducta, bien jurídico, eventual resultado y gravedad deben ser identificables con suficiente claridad.\par
            La indeterminación grave puede afectar la reserva de ley y el conocimiento previo de la prohibición.\par





          Lex praevia\par

            Consagra la irretroactividad in malam partem y la predictibilidad: no puede penarse lo que la ley no conminó con anterioridad\par
            suficiente, salvo regímenes de favorabilidad expresos.\par





          Lex stricta\par

            Impone interpretación restrictiva en favor del imputado y límites severos a la analogía que agrave la situación del acusado.\par
            Es el corolario dogmático de la seguridad jurídica en materia penal.\par









      5. Cuadro operativo (litigio)\par



          DimensiónPregunta de trabajoSalida procesal frecuente\par




          Lex certa¿Puedo describir la conducta típica sin rellenar vacíos discrecionales ilegítimos?Acción de inconstitucionalidad; nulidad; excepciones según etapa\par

          Lex praevia¿La conducta estaba conminada con ley válida y anterior?Excepción de legalidad; absolución; revisión de temporalidad\par

          Lex stricta¿La interpretación aplicada es la más favorable al imputado compatible con el texto?Revisión en casación; anulación por error jurídico\par








      6. Autocomprobación\par

      Ordene mentalmente los cuatro bloques del mapa (rectores → dignidad → bloque → legalidad). Luego pulse:\par

      Ver criterio de cierre\par





      Fuentes: Temario L1 (Conceptos). Bernal, Derecho penal general, Cap. 4: pp. 251–266 (dignidad), 266–285 (bloque),\par
      285–299 (principios de sanciones), 341–366 (legalidad). Capítulo 4 completo: pp. 251–418.\par


    Siguiente: Explicación →\par

\section{Explicaci\'on}
\noindent\textit{Archivo fuente:} \texttt{02-explicacion.html}

L1 — Explicación\par





    Elemento 2 — Explicación\par
    Integración constitucional y precedentes orientadores\par

    Introducción breve: cómo el bloque de constitucionalidad y la Corte Constitucional moldean la interpretación penal.\par


      Este elemento explica la cadena argumentativa que va del texto constitucional y los tratados al juzgamiento penal concreto.\par
      Los precedentes seleccionados ilustran criterios reiterados sobre integración normativa, garantías del imputado y límites a la\par
      expansión judicial del tipo. La tabla resume enseñanzas operativas para el litigante; el detalle doctrinal completo está en Bernal, Cap. 4.\par



      1. Bloque de constitucionalidad en acción\par


        El bloque lato sensu obliga al intérprete a situar el tipo penal dentro de un sistema de fuentes donde conviven Constitución,\par
        leyes estatutarias, tratados incorporados y estándares de derechos humanos. El bloque stricto sensu refuerza la supremacía:\par
        normas inferiores incompatibles quedan sin aplicación o son objeto de control concentrado.\par



        1.1 Efectos en la sala penal\par

          El juez penal no “crea” el bloque, pero debe aplicar interpretaciones conformes que eviten colisiones con derechos\par
          fundamentales. Eso condiciona la lectura de expresiones indeterminadas del tipo, la ponderación de circunstancias y la revisión de\par
          medidas cautelares personales.\par




        1.2 Relación con lex stricta\par

          Cuanto más amplio es el estándar de protección constitucional, más exigente resulta la justificación de una interpretación extensiva\par
          del tipo. La analogía in malam partem queda vedada; la analogía interpretativa favorable al imputado puede ser admisible solo\par
          en los límites que el ordenamiento y la jurisprudencia de unificación permitan.\par






      2. Matriz de precedentes (síntesis para litigio)\par


        Cada fila resume la línea jurisprudencial útil para el penalista, sin sustituir la lectura del fallo íntegro ni de la\par
        guía de estudio (Bernal, Cap. 4, apartados sobre control constitucional y bloque).\par





            Precedente\par
            Tema central\par
            Criterio clave para el penalista\par
            Efecto práctico\par






            C-225 de 1995\par
            Integración del bloque de constitucionalidad y tratados sobre derechos humanos\par
            Los instrumentos internacionales ratificados que consagren derechos humanos forman parte del acervo interpretativo que informa la aplicación del ordenamiento, en la medida en que el sistema de fuentes los incorpore según la Constitución y la ley.\par
            Fortalece argumentos de inconstitucionalidad o de interpretación conforme cuando una norma o práctica penal colisiona con estándares convencionales ya incorporados al bloque.\par



            C-582 de 1999\par
            Garantías del imputado y estándares de debido proceso en el proceso penal\par
            La afectación grave e irremediable de garantías sustanciales puede proyectarse en nulidad o en líneas defensivas que exigen reparación procesal efectiva, no meramente formal.\par
            Útil para articular excepciones, incidentes y recursos cuando hay vulneración estructural del derecho de defensa o del contradictorio.\par



            C-146 de 2021\par
            Inhabilidades para ser elegido (Ley 136 de 1994) y articulación del bloque de constitucionalidad con el Derecho Internacional de los Derechos Humanos\par
            La Corte examina límites razonables al derecho político de elegir y ser elegido y, en paralelo, precisa el papel de la jurisprudencia de la Corte IDH: útil como criterio hermenéutico al interpretar el bloque, sin convertir a la Corte Constitucional en “juez de convencionalidad” en el sentido de confrontación abstracta ley-tratado que desplace la supremacía constitucional.\par
            Para el penalista: al invocar tratados o fallos IDH ante juez ordinario, la argumentación debe articularse en términos de interpretación conforme y armonización con la Constitución, evitando trasplantes automáticos de estándares sin anclaje en el parámetro interno.\par







        Para citas textuales en escritos judiciales o académicos, consulte siempre el texto oficial del fallo en la Corte Constitucional y cruce con Bernal, Cap. 4 (pp. 266–285, 341–366).\par





      3. Secuencia de razonamiento (de principio a subsumición)\par



          Paso 1 — Marco\par
          Identificar qué derecho fundamental o principio rector está en juego (dignidad, legalidad, debido proceso).\par




          Paso 2 — Estándar\par
          Fijar el estándar de protección: constitucional, convencional o de bloque integrado.\par




          Paso 3 — Colisión\par
          Describir cómo la norma, la interpretación o el acto procesal vulneran ese estándar.\par




          Paso 4 — Remedio\par
          Elegir la vía idónea: interpretación conforme, excepción, incidente, recurso de casación o control abstracto, según competencia y etapa.\par








      4. Contraste: expansión del tipo frente a reserva de ley\par



          PosturaArgumento típicoContrapeso constitucional\par





            Interpretación extensiva acusatoria\par
            Ampliar el significado de un verbo o resultado para cubrir conductas “equivalentes”\par
            Lex stricta y tipicidad: la expansión sin respaldo legal claro erosiona lex certa y el conocimiento previo\par



            Interpretación restrictiva defensiva\par
            Cerrar el tipo al sentido más favorable al imputado compatible con el texto y la finalidad de la norma\par
            Coherente con el principio de inocencia y con la función garantista del derecho penal\par









      5. Autocomprobación\par

      ¿Cuál precedente refuerza con mayor nitidez la integración de tratados DDHH en el parámetro interpretativo colombiano?\par

      C-225 de 1995\par
      C-582 de 1999\par





      Fuentes: Temario L1 (Explicación). Bernal, Cap. 4: bloque de constitucionalidad (pp. 266–285), legalidad (pp. 341–366).\par
      Fallos C-225/95, C-582/99, C-146/21 (consulta oficial Corte Constitucional).\par


    Siguiente: Exploración →\par

\section{Exploraci\'on}
\noindent\textit{Archivo fuente:} \texttt{03-exploracion.html}

L1 — Exploración\par





    Elemento 3 — Exploración\par
    Dilemas: legalidad estricta, bloque y estándares internacionales\par

    Introducción breve: explorar posturas sin sustituir el análisis del caso concreto.\par


      La exploración confronta tres tensiones recurrentes: (1) tipicidad estricta frente a integración de categorías de Derecho Internacional;\par
      (2) seguridad jurídica frente a exigencias de no impunidad en graves violaciones; (3) coordinación procesal\par
      (ne bis in idem, competencia) frente a subsunción sustantiva única. Cada microcaso pide una primera línea argumentativa y\par
      devuelve criterios de autocorrección.\par



      Microcaso 1 — Tipo interno vs. categoría internacional\par


        Hechos sintéticos: investigación por homicidio agravado; la Fiscalía insinúa encuadre como crimen de lesa humanidad ante posible participación\par
        en patrón sistemático. La defensa invoca lex praevia y ausencia de tipificación expresa de “crimen de lesa humanidad” en el momento relevante.\par


       A) Cerrar el debate en el tipo penal interno vigente; no admitir categoría internacional en primera instancia.\par
       B) Admitir integración inmediata de categoría internacional aunque el tipo interno sea distinto.\par
       C) Separar: primero competencia y temporalidad; luego, si procede, subsidiariamente armonizar con DIH.\par
      Retroalimentación\par





      Microcaso 2 — Interpretación extensiva del verbo típico\par


        El tipo usa un verbo de amplia extensión semántica en la práctica forense. El juez tiende a equiparar dos conductas por “equivalencia social”\par
        sin norma aclaratoria. La defensa alega lex stricta.\par


       A) Apoyar la equiparación si protege el bien jurídico.\par
       B) Exigir que toda ampliación relevante provenga del legislador o de criterios de unificación vinculantes.\par
       C) Proponer interpretación teleológica acusatoria sin límites textuales.\par
      Retroalimentación\par





      Microcaso 3 — Tratado y bloque ante juez penal\par


        La defensa cita un fallo de la Corte IDH y pide inaplicar una norma interna. El juez penal duda de su alcance tras la línea de C-146/21 sobre\par
        el papel de la jurisprudencia interamericana y la supremacía constitucional.\par


       A) Solicitar inaplicación automática de la norma interna.\par
       B) Articular interpretación conforme: mostrar colisión, vía procesal idónea y estándar del bloque interno.\par
       C) Ignorar el tratado por ser “externo” al proceso penal.\par
      Retroalimentación\par





      Simulador global de postura (macro)\par

      Elija la línea que defendería como estrategia general (autoevaluación):\par

       A) Prioridad absoluta a legalidad estricta y límites interpretativos.\par
       B) Prioridad a estándares internacionales cuando el orden interno genera impunidad frente a graves violaciones.\par
       C) Estrategia dual: competencia, temporalidad y ne bis in idem primero; integración subsidiaria.\par
      Ver guía de argumentos\par





      Tabla de tensión (refuerzo)\par



          VariableRiesgo si se privilegia solo un poloContrapeso dogmático\par




          Legalidad estrictaFormalismo que obstruye protección efectiva de víctimas en macro-crimenReserva de ley, tipicidad, inocencia\par

          Integración internacionalInseguridad jurídica y choque con texto interno sin canal procesalBloque, armonización, vías expresas\par

          Proporcionalidad de la intervenciónSobrepunitivismo o subprotección según el sesgo institucionalDignidad humana, fines constitucionales de la pena\par








      Fuentes: Temario L1 (Exploración). Bernal, Cap. 4 (principios rectores, bloque, legalidad). Sentencia C-146/21 (papel hermenéutico de Corte IDH y bloque).\par


    Siguiente: Contexto →\par

\section{Contexto}
\noindent\textit{Archivo fuente:} \texttt{04-contexto.html}

L1 — Contexto\par





    Elemento 4 — Contexto\par
    Ambigüedad legislativa y estándar de control\par

    Introducción breve: contextualizar por qué la vaguedad del lenguaje penal dispara conflictos de constitucionalidad.\par


      El legislador usa términos evaluativos (“grave”, “abusivo”, “arbitrario”) o técnicos abiertos. El juez debe concretarlos sin convertirse en\par
      legislador positivo. El contexto constitucional exige: (i) identificar el grado de indeterminación; (ii) aplicar criterios de\par
      lex certa y proporcionalidad; (iii) preferir interpretaciones que preserven la predictibilidad sin vaciar la protección del bien jurídico.\par



      1. Tipologías de ambigüedad (laboratorio conceptual)\par



          TipoEjemplo genéricoProblema dogmáticoPregunta de control\par




          SemánticaVerbo polisémico en el tipoRiesgo de expansión judicial¿El ciudadano común puede prever la conducta prohibida?\par

          EvaluativaGravedad, proporcionalidad en el tipoDelegación implícita al juzgador¿Existen parámetros legales o jurisprudencia unificada que acoten?\par

          SistémicaColisión entre tipos o con normas especialesDoble valoración o lagunas¿Cuál es la especialidad relativa y el criterio de consumición?\par

          TecnológicaConductas digitales no previstas literalmenteAnalogía y tipicidad¿La conducta cae por interpretación restrictiva o requiere reforma legislativa?\par








      2. Criterio de gravedad para el control\par


        No toda imprecisión invalida el tipo. La doctrina constitucional y penal distingue entre indeterminación tolerable\par
        (núcleo claro del tipo, márgenes de juicio previstos) e indeterminación inconstitucional (núcleo inidentificable o\par
        discrecionalidad sin ley).\par




          Señales de tolerabilidad\par


-- El núcleo de la conducta está delimitado por el texto y la práctica consolidada.\par


-- La evaluación remite a criterios objetivos o a listas legales.\par


-- La defensa puede anticipar la línea acusatoria con trabajo de investigación razonable.\par





          Señales de riesgo alto\par


-- El tipo se rellena solo con moral social o conveniencia procesal.\par


-- La conducta punible depende de opinión del funcionario sin estándar escrito.\par


-- La misma acción puede ser típica o atípica según el juzgado sin reglas comunes.\par









      3. Ejemplo analítico (sintético)\par


        Un tipo penal sanciona “trato degradante” sin definición legal ni remisión a catálogo. La defensa invoca lex certa. La Fiscalía\par
        propone una lista casuística extraída de la práctica internacional. El juez debe: verificar si el núcleo es al menos previsible por\par
        referencia a fuentes válidas del ordenamiento; si no, la vía puede ser inconstitucionalidad selectiva o interpretación restrictiva que excluya\par
        la conducta concreta.\par



        Clave de lectura (Bernal)\par

          Relacione este ejemplo con el bloque sobre legalidad estricta y principios rectores (Bernal, Cap. 4,\par
          especialmente la sección de legalidad, pp. 341–366): la solución no es “más punitivismo” sino coherencia entre reserva de ley y estándar\par
          interpretativo.\par






      4. Autocomprobación\par

      ¿Qué tipo de ambigüedad es más sensible para lex stricta?\par

      Evaluativa\par
      Semántica del verbo típico\par





      Fuentes: Temario L1 (Contexto). Bernal, Cap. 4: legalidad (pp. 341–366), principios y bloque (pp. 251–299, 266–285).\par


    Siguiente: Aplicación →\par

\section{Aplicaci\'on}
\noindent\textit{Archivo fuente:} \texttt{05-aplicacion.html}

L1 — Aplicación\par





    Elemento 5 — Aplicación\par
    Clínica: recurso extraordinario y error jurídico en clave constitucional\par

    Introducción breve: traducir principios rectores y legalidad en un esquema de impugnación estructurado.\par


      Este elemento aplica la teoría de la Lección 1 a un supuesto de segunda instancia (o casación, según el proceso): se busca\par
      demostrar error jurídico en la interpretación del tipo o en la ponderación de derechos fundamentales, articulando bloque de\par
      constitucionalidad y estándar de legalidad estricta. La plantilla siguiente está redactada en tono forense definitivo, lista para adaptar a hechos probados.\par



      1. Hechos relevantes y decisión impugnada (extracto)\par


        La sentencia de instancia interpretó extensivamente el verbo típico para incluir conducta B, aunque el relato probatorio\par
        describe conducta A, asimilable solo por vía de “equivalencia social”. No existe norma aclaratoria ni línea unificada\par
        vinculante que autorice la equiparación. La defensa sostiene vulneración de lex certa y lex stricta, con proyección al\par
        principio de inocencia y al debido proceso.\par





      2. Tesis jurídicas (numeradas)\par



-- La interpretación extensiva del tipo sin soporte legal claro contraviene el mandato de legalidad estricta (Constitución Política; desarrollo dogmático en Bernal, Cap. 4, sección de legalidad).\par


-- La equiparación A ≈ B importa analogía in malam partem prohibida, al agraviar la situación del imputado sin base textual suficiente.\par


-- El estándar del bloque de constitucionalidad exige interpretación conforme que preserve derechos fundamentales; la sentencia impugnada no explicitó el test de certeza ni el estándar restrictivo aplicable.\par





      3. Fundamentos en bloque y precedentes (citación de trabajo)\par


        En el escrito definitivo se citará el texto legal pertinente y, si aplica, C-225/95 (integración del bloque y tratados),\par
        C-582/99 (garantías del imputado) y la precisión sobre el valor hermenéutico de la jurisprudencia interamericana en\par
        C-146/21, para mostrar que la integración no autoriza saltos interpretativos sin encadenamiento al parámetro interno.\par




          ArgumentoSoporteConclusión parcial\par




          Certeza del tipoLex certa; doctrina de reserva de leyLa conducta punible debe ser identificable sin rellenar el tipo con criterios extralegales discrecionales\par

          Límite interpretativoLex strictaCorresponde la interpretación más favorable al imputado compatible con el texto y la finalidad constitucional de la norma\par

          Debido procesoArts. 28 y 29 C.P.; línea C-582/99La imputación debe ser precisa y contradictoria; la ampliación del tipo en la sentencia afecta la defensa técnica\par








      4. Petitorio\par



-- Revocar la sentencia impugnada en el capítulo de tipicidad objetiva y, subsidiariamente, en la valoración de la intervención sobre libertad si a ello hubiere lugar.\par


-- Reenviar para nuevo fallo con obligación de motivar bajo estándar de legalidad estricta y bloque de constitucionalidad.\par


-- Imponer costas conforme a ley.\par





      5. Campos para adaptación (su caso)\par

      Sustituya los marcadores por datos de su expediente.\par



      Verificar que completó ambos campos\par





      Fuentes: Temario L1 (Aplicación). Bernal, Cap. 4: principios rectores, bloque, legalidad (pp. 251–418, con énfasis 341–366).\par


    Siguiente: Prospectiva →\par

\section{Prospectiva}
\noindent\textit{Archivo fuente:} \texttt{06-prospectiva.html}

L1 — Prospectiva\par





    Elemento 6 — Prospectiva\par
    Escenarios futuros y priorización estratégica\par

    Introducción breve: anticipar reformas, litigiosidad y riesgos de política criminal.\par


      La prospectiva no predice el futuro: ordena variables (legislación, jurisprudencia constitucional, estándares internacionales,\par
      tecnología) y asigna probabilidad de ocurrencia y impacto en el litigio penal. La tabla combina escenarios\par
      con mitigaciones concretas para abogacía y academia.\par



      1. Matriz de escenarios\par




            Escenario\par
            Disparadores\par
            Efecto probable en penal\par
            Prioridad táctica\par






            Endurecimiento punitivo legislativo\par
            Crisis de seguridad, presión electoral\par
            Tipos ampliados, penas máximas alzas; más tensiones con lex certa\par
            Alta: preparar líneas de inconstitucionalidad y defensa de reserva de ley\par



            Mayor control constitucional concentrado\par
            Fallos de inexequibilidad en materia penal o procesal\par
            Reordenamiento de tipos o pruebas; inseguridad transitoria\par
            Media-alta: monitoreo de relatoría y ajuste de escritos modelo\par



            Integración hermenéutica de DIH sin nueva codificación\par
            Sentencias IDH + debates académicos\par
            Más argumentos de interpretación conforme en sala penal\par
            Media: dominar C-225/95 y límites de C-146/21\par



            Criminalización de conductas digitales\par
            Nuevas tecnologías, tipos “omnibus”\par
            Conflictos de tipicidad, prueba electrónica, proporcionalidad\par
            Alta: pericia y tipicidad estricta\par



            Desjudicialización o negociación reforzada\par
            Sobrecarga judicial, políticas restaurativas\par
            Cambio de incentivos; menos litigio en algunos delitos menores\par
            Baja-media según especialidad; seguir desarrollo legislativo\par









      2. Riesgos y mitigaciones\par



          Riesgo institucional\par



-- Inflación tipológica sin política criminal clara → mitigar con estudios de impacto y control abstracto.\par


-- Fragmentación jurisprudencial entre salas → mitigar con recursos de unificación y citas sistemáticas.\par





          Riesgo para el litigante\par



-- Sobrerreacción defensiva (demasiados motivos dilatorios) → priorizar ataques constitucionalmente densos.\par


-- Subutilización del bloque → formación continua en tratados y estándar C-146/21.\par









      3. Priorización (autoevaluación)\par

      Ordene mentalmente: ¿qué escenario afecta más su práctica en los próximos 24 meses?\par

      Endurecimiento legislativo\par
      Digital / prueba electrónica\par





      Fuentes: Temario L1 (Prospectiva). Bernal, Cap. 4 (visión sistémica de principios y límites al ius puniendi, pp. 251–418).\par


    Siguiente: Ejercicios →\par

\section{Ejercicios}
\noindent\textit{Archivo fuente:} \texttt{07-ejercicios.html}

L1 — Ejercicios\par





    Elemento 7 — Ejercicios\par
    Ejercicios guiados con justificación modelo\par

    Introducción breve: seis ítems de distinta dificultad; cada uno incluye respuesta modelo al revelar.\par


      Los ejercicios recorren definición, distinción, aplicación y síntesis.\par
      Use primero su criterio; luego pulse “Ver modelo” para contrastar con la pauta académica alineada a Bernal, Cap. 4.\par



      Ej. 1 — Definición operativa\par

      En una frase técnica, defina bloque de constitucionalidad lato sensu.\par

      Ver modelo\par





      Ej. 2 — Contraste\par

      Enumere dos diferencias entre lex certa y lex stricta.\par

      Ver modelo\par





      Ej. 3 — Subsumición breve\par

      Un juez amplía el tipo por analogía que agrava. ¿Qué principio o principios invoca usted en defensa y en qué orden?\par

      Ver modelo\par





      Ej. 4 — Precedentes\par

      ¿Qué aporta C-225/95 al argumento penal que invoca un tratado de derechos humanos ratificado?\par

      Ver modelo\par





      Ej. 5 — C-146/21\par

      Explique en dos líneas por qué C-146/21 no autoriza al abogado penal a pedir “inaplicación automática” de la ley interna citando solo un fallo de la Corte IDH.\par

      Ver modelo\par





      Ej. 6 — Síntesis\par

      Elabore un esquema de cuatro pasos del razonamiento constitucional penal (del principio a la vía procesal).\par

      Ver modelo\par





      Fuentes: Temario L1 (Ejercicios). Bernal, Cap. 4 (pp. 251–418).\par


    Siguiente: Quiz →\par

\section{Quiz}
\noindent\textit{Archivo fuente:} \texttt{08-quiz.html}

L1 — Quiz\par





    Elemento 8 — Quiz\par
    Evaluación L1: principios rectores y bloque\par

    12 preguntas. Al enviar verá puntuación y retroalimentación por ítem fallado.\par


      Cada pregunta tiene una sola respuesta correcta. El quiz está alineado a la Lección 1 (Bernal, Cap. 4, pp. 251–418) y a los elementos previos del curso.\par




        1. El bloque de constitucionalidad lato sensu se distingue porque:\par

         A) Solo incluye la Constitución Política.\par
         B) Integra un conjunto amplio de fuentes que informan la interpretación conforme del ordenamiento.\par
         C) Excluye tratados internacionales ratificados.\par



        2. Lex praevia se asocia principalmente con:\par

         A) Interpretación restrictiva del tipo.\par
         B) Irretroactividad de la ley penal en perjuicio del imputado y predictibilidad.\par
         C) Claridad absoluta del lenguaje del juez.\par



        3. En la línea de C-225/95, los tratados sobre DDHH ratificados cumplen rol relevante como:\par

         A) Fuentes integrantes del parámetro interpretativo del bloque, según el sistema de fuentes.\par
         B) Normas siempre superiores al código penal sin más.\par
         C) Opinión consultiva sin efectos en juzgamiento.\par



        4. Lex stricta exige:\par

         A) Interpretación restrictiva en favor del imputado y límites severos a analogías que agraven.\par
         B) Que el juez elija la pena sin reglas.\par
         C) Solo aplicar costumbre penal.\par



        5. La dignidad humana en el derecho penal cumple función:\par

         A) Meramente declarativa sin incidencia en la pena.\par
         B) Teleológica: humanización y proporcionalidad de la intervención punitiva.\par
         C) Excluir el debido proceso.\par



        6. La analogía in malam partem:\par

         A) Está vedada en el sistema de legalidad estricta.\par
         B) Es el criterio preferente del juez penal.\par
         C) Solo aplica a multas administrativas.\par



        7. C-582/99 es especialmente útil para argumentar:\par

         A) Garantías del imputado y estándares de debido proceso en el proceso penal.\par
         B) Impuestos municipales.\par
         C) Competencia civil exclusivamente.\par



        8. Tras la línea de C-146/21, las sentencias de la Corte IDH:\par

         A) Sirven como criterios hermenéuticos al interpretar el bloque, sin automatizar la inaplicación de la ley interna.\par
         B) Obligan siempre a inaplicar cualquier norma interna sin análisis.\par
         C) Carecen de relevancia para el penalista.\par



        9. La claridad y predeterminación del contenido del tipo se vinculan sobre todo con:\par

         A) Lex certa.\par
         B) Lex stricta únicamente.\par
         C) Costumbre mercantil.\par



        10. El bloque stricto sensu refuerza:\par

         A) La supremacía del núcleo constitucional y mecanismos de control frente a normas incompatibles.\par
         B) La exclusión total de tratados.\par
         C) La discrecionalidad ilimitada del fiscal.\par



        11. Una expansión judicial del verbo típico que agrava al imputado sin base legal clara vulnera prioritariamente:\par

         A) Solo el principio de inocencia, nunca la legalidad.\par
         B) Lex stricta y, en conexión, lex certa y legalidad material.\par
         C) Únicamente la caducidad.\par



        12. En la secuencia de razonamiento constitucional penal del curso, el primer paso es:\par

         A) Identificar el remedio procesal sin más.\par
         B) Identificar qué derecho fundamental o principio rector está en juego.\par
         C) Pedir conciliación fiscal.\par


      Enviar y ver resultados\par



      Resultado\par










      Fuentes: Temario L1 (Quiz). Bernal, Cap. 4. Sentencias C-225/95, C-582/99, C-146/21.\par


    Volver al inicio L1\par

\end{document}

